В первой главе маркетинговый сайт рассматривается не как набор отдельных страниц, а как
предметная область, в которой эксплуатационная ценность системы определяется скоростью,
управляемостью и предсказуемостью контентных изменений. Поэтому аналитика выстроена в
последовательности ``особенности предметной области -> архитектурный выбор -> требования
к системе'': сначала выделяются свойства маркетинговой витрины как объекта автоматизации,
затем оцениваются подходы к управлению контентом и после этого формируется набор
требований, который используется как основание для проектных решений второй главы. В
рамках ВКР такая логика принципиальна, поскольку разрабатывается не просто публичный
frontend, а \texttt{CMS-first} система управления контентом для корпоративной
маркетинговой платформы.

\section{Особенности маркетинговых сайтов как объекта автоматизации}

\subsection{Роль маркетингового сайта в цифровых коммуникациях компании}

Маркетинговый сайт в современной организации выполняет роль публичной цифровой витрины,
через которую компания представляет продукты и услуги, формирует профессиональную
репутацию, публикует экспертные материалы, демонстрирует выполненные проекты и поддерживает
коммуникацию с потенциальными клиентами и соискателями. В отличие от простого
информационного сайта, он должен быстро адаптироваться к новым маркетинговым кампаниям,
изменениям позиционирования, запуску новых направлений и потребности в регулярном
обновлении содержимого~\cite{torgunakova2023marketing-website,voroniuk2020retail-cms-needs}.

Однако для инженерного анализа важен не сам факт наличия таких материалов, а режим их
изменения. В маркетинговом контуре содержание сайта обновляется не эпизодически, а в
ответ на кампании, запуск новых предложений, публикацию экспертных материалов, изменения
позиционирования и кадровые сценарии. Следовательно, система должна поддерживать частые
редакторские операции как нормальный режим работы, а не как исключение. Это смещает
фокус проектирования с разовой разработки шаблонов на организацию управляемого
контентного контура, в котором типовые изменения не требуют вмешательства в код
публичной витрины.

Из этого вытекает и второй важный вывод: ценность маркетингового сайта определяется не
только качеством визуального представления, но и тем, насколько быстро и предсказуемо
он позволяет провести материал через цикл ``подготовка -> проверка -> публикация''.
Если типовые контентные изменения зависят от разработчика, возникает организационное
узкое место: маркетинговая команда не может оперативно обновлять страницы, публикации и
метаданные, а сопровождение сайта начинает смешиваться с задачами развития frontend.
Следовательно, объектом автоматизации становится не только отображение контента, но и
сам редакторский процесс, включая его роли, точки контроля и публикационный контур.

Для корпоративного маркетингового сайта также существенны требования к поисковой
видимости, локализации и управляемой публикации. Контент должен быть пригоден для
индексации, сопровождаться метаданными, поддерживать локализованные версии и
публиковаться в воспроизводимом контуре, не зависящем от ручных действий в кодовой
базе. Эти свойства делают сайт не просто набором маршрутов, а контентной системой со
своими пользовательскими ролями, моделями данных и эксплуатационными сценариями. Для
русскоязычных работ по
информационной архитектуре и корпоративным сайтам также характерен акцент на устойчивой
навигации, структурной согласованности контента и централизованном управлении
информационным слоем~\cite{belyaev2009website-navigation,mikhailova2016content-navigation,mashina2022centralized-cms-knowledge}.

В таблице~\ref{tab:marketing-domain-engineering-effects} сведены ключевые свойства
рассматриваемой предметной области и их прямые инженерные последствия для проектируемой
системы.

\begin{table}[htbp]
  \centering
  \small
  \caption{Связь особенностей маркетингового сайта с инженерными требованиями к системе}
  \label{tab:marketing-domain-engineering-effects}
  \begin{tabular}{|p{0.34\linewidth}|p{0.56\linewidth}|}
    \hline
    \textbf{Особенность предметной области} & \textbf{Инженерное следствие для системы} \\
    \hline
    Частые изменения контента в ответ на кампании, публикации и изменения позиционирования &
    Типовые обновления должны выполняться через управляемые контентные сущности и
    редакторский интерфейс, без ручного изменения frontend-кода. \\
    \hline
    Необходимость независимого редакторского цикла &
    Система должна разделять роли редактирования и разработки, поддерживать черновой
    режим, проверку материала и публикацию как самостоятельный контур. \\
    \hline
    Мультиязычность публичного контента &
    Локализация должна проектироваться на уровне данных, маршрутов и редакторских
    сценариев, а не сводиться к разрозненным текстовым вставкам в шаблонах. \\
    \hline
    Управление \texttt{SEO} и метаданными &
    Метаданные должны быть частью контентной модели, чтобы редактор мог управлять
    индексируемостью, заголовками, описаниями и \texttt{Open Graph} без участия
    разработчика. \\
    \hline
    Необходимость предварительного просмотра до публикации &
    Требуется отдельный сценарий \texttt{preview}, позволяющий проверить будущую
    публичную версию до вывода материала в основной контур. \\
    \hline
    Требование воспроизводимой схемы публикации &
    Публикация должна рассматриваться как формализованный процесс обновления витрины, а
    не как набор ручных операций в кодовой базе и инфраструктуре. \\
    \hline
  \end{tabular}
\end{table}

\subsection{Контентные сущности и редакторские сценарии}

Для рассматриваемой предметной области характерен ограниченный, но содержательно важный
набор сущностей. К ним относятся произвольные маркетинговые страницы, статьи, проекты,
вакансии, отклики на вакансии, а также глобальные данные сайта. Каждая из этих сущностей
имеет собственный жизненный цикл и собственные требования к редактированию. Например,
маркетинговая страница должна собираться из переиспользуемых блоков, статья требует
структурированного текста, обложки и авторства, проект должен показывать портфолио
выполненных работ, а вакансия должна не только отображаться на витрине, но и быть связана
со сценарием отправки отклика.

Инженерно значимо то, что перечисленные сущности не сводятся к одному шаблону
``страница с текстом''. Они требуют разных структур данных, разных правил публикации и
разных связей между редакторским и публичным контуром. Если система хранит такие
материалы как слабо структурированный набор полей, то редакторский процесс быстро
становится зависимым от частных доработок интерфейса и ручной правки представления.
Следовательно, уже на уровне предметной области возникает требование к явной
контентной модели, в которой страницы, статьи, проекты, вакансии и глобальные данные
оформлены как самостоятельные сущности со своими атрибутами и жизненным циклом.

Отсюда вытекает основной редакторский сценарий: контент-менеджер или маркетолог должен
создать или изменить материал, заполнить структурированные поля, при необходимости
подготовить локализованную версию, проверить публичное представление и опубликовать
результат без ручного изменения frontend-кода. В этой логике \texttt{SEO}-поля,
метаданные \texttt{Open Graph}, локализованные маршруты и предварительный просмотр
перестают быть второстепенными опциями интерфейса и становятся обязательными элементами
архитектуры редакторского контура.

Таким образом, маркетинговый сайт представляет собой систему, в которой ценность создается
не только конечной HTML-страницей, но и качеством управления ее содержимым. Это отличает
рассматриваемую предметную область от проектов, где достаточно один раз сверстать
несколько шаблонов и редко возвращаться к их содержимому. Для ВКР по разработке системы
управления контентом такое различие принципиально: требуется проектировать именно
управляемую контентную платформу, а не только витринный интерфейс. Именно поэтому
следующим шагом анализа становится сравнение архитектурных подходов к управлению
контентом: необходимо определить, какая модель \texttt{CMS} лучше соответствует
требованиям независимого редакторского цикла, мультиязычности, управляемого
\texttt{SEO}, \texttt{preview} и воспроизводимой публикации.

\section{Подходы к управлению контентом: традиционные CMS и headless CMS}

\subsection{Ограничения традиционного CMS-подхода}

Традиционная \texttt{CMS} обычно объединяет в одном контуре хранение данных, административный
интерфейс, шаблоны отображения и механизм вывода страницы пользователю. Такая архитектура
хорошо подходит для типовых сайтов, в которых структура представления относительно стабильна,
а публичный слой тесно связан с самой системой управления контентом. Для небольших задач это
снижает порог входа и ускоряет старт разработки, поскольку многие возможности доступны из
коробки~\cite{yesikov2013cms-comparative-analysis}.
Подобная логика характерна и для \texttt{CMS}-ориентированных framework-подходов,
автоматизирующих создание типовых веб-сайтов и CRUD-контуров, где приоритет отдается
скорости запуска и унификации административной части~\cite{bandirmali2018mtcmf}.

Для задачи маркетингового сайта принципиально, что преимущества такого подхода относятся
прежде всего к стартовой простоте, а не к гибкости архитектуры. Публичная витрина в этом
случае зависит от шаблонного и runtime-слоя самой \texttt{CMS}, поэтому развитие frontend
обычно происходит внутри ограничений конкретной темы, движка и способа рендеринга. Это
усложняет ситуацию, когда требуется независимо изменять структуру маркетинговых страниц,
переиспользовать одни и те же контентные сущности в разных разделах или выносить часть
публичного слоя в отдельный технологический контур.

Интеграция редактирования и рендеринга действительно упрощает базовый \texttt{preview} и
публикацию: редактор работает в той же системе, которая формирует страницу. Однако для
рассматриваемой предметной области этого недостаточно. Если локализация, \texttt{SEO},
метаданные и публикационный процесс завязаны на тему и внутренний runtime \texttt{CMS},
то они начинают зависеть не столько от контентной модели, сколько от частных доработок
слоя представления. В результате традиционный подход не исключает нужные функции, но
делает их вторичными по отношению к связанному шаблонному контуру, тогда как в ВКР
требуется архитектура, где редакторский и публичный слои могут развиваться раздельно.

\subsection{Преимущества headless CMS для маркетинговой платформы}

\texttt{Headless CMS} исходит из другой модели: система управления отвечает за хранение
и редактирование структурированных данных, а отображение контента переносится в отдельный
frontend-слой. Такой подход естественным образом соответствует \texttt{API-first} логике,
в которой \texttt{CMS} предоставляет данные через программный интерфейс, а витрина или
другой потребитель решает, как именно эти данные представить пользователю. Для
\texttt{CMS}-ориентированной разработки это означает более явное разделение контентной
модели, правил публикации и публичного интерфейса, что хорошо согласуется с современными
model-driven и platform-independent подходами в домене \texttt{CMS}~\cite{priefer2021cms-mdd-domain,gkantouna2019cms-oriented-modeling,filipe2016xis-cms}.

Для маркетингового сайта эта модель важна не сама по себе, а как способ перевести
ключевые требования предметной области в архитектурные свойства системы. Когда frontend
не встроен в \texttt{CMS}, контентная модель может проектироваться относительно сущностей
и редакторских сценариев, а не относительно одного набора шаблонов. Это упрощает
переиспользование контента, позволяет строить locale-aware маршрутизацию и дает основу
для отдельного публикационного контура, в котором обновление данных не требует ручной
правки публичной витрины~\cite{santahuhta2022headless-cms,matjuhhin2026api-driven-cms,holcik2025headless-cloud}.

При этом \texttt{headless} подход не снимает инженерные задачи автоматически. В частности,
\texttt{preview}, управление метаданными и публикация должны проектироваться явно, так как
\texttt{CMS} не рендерит итоговую страницу самостоятельно. Это усложняет начальную
интеграцию, но одновременно делает публикационный контур явным и проверяемым: редакторский
слой, \texttt{API}, frontend-сборка и публичная выдача перестают смешиваться в одном
runtime-сценарии. Для мультиязычных контентных платформ это особенно существенно, поскольку
локаль влияет не только на тексты интерфейса, но и на маршруты, метаданные и доступность
данных в разных языковых версиях~\cite{santahuhta2022headless-cms,oleary2008multilingual-kms,gracia2012multilingual-web-data,solovev2023multilingual-web-app}.

Для формализации выбора архитектурного подхода под задачу маркетинговой витрины
целесообразно сопоставить оба варианта по критериям, которые прямо следуют из предметной
области и целевого результата ВКР (таблица~\ref{tab:cms-vs-headless-marketing}).

\begin{table}[htbp]
  \centering
  \footnotesize
  \caption{Сравнение традиционной CMS и headless CMS по критериям маркетингового сайта}
  \label{tab:cms-vs-headless-marketing}
  \setlength{\tabcolsep}{4pt}
  \begin{tabular}{|p{0.16\linewidth}|p{0.20\linewidth}|p{0.20\linewidth}|p{0.26\linewidth}|}
    \hline
    \textbf{Критерий} & \textbf{Традиционная CMS} & \textbf{Headless CMS} & \textbf{Значение для задачи ВКР} \\
    \hline
    Зависимость frontend-слоя от \texttt{CMS} &
    Шаблоны и рендеринг обычно встроены в \texttt{CMS}; развитие витрины зависит от ее theme- и runtime-слоя. &
    Публичный frontend развивается отдельно и использует \texttt{CMS} как источник данных по \texttt{API}. &
    Для независимого развития витрины и редакторского контура предпочтительно слабое связывание представления с \texttt{CMS}. \\
    \hline
    Пригодность к мультиязычности &
    Локализация часто расширяется на уровне шаблонов, плагинов и отдельных страниц. &
    Локаль можно проектировать на уровне данных, \texttt{API}, маршрутов и метаданных как единый контур. &
    Для обязательной схемы \texttt{ru/en} требуется архитектурная, а не точечная локализация. \\
    \hline
    Повторное использование контента &
    Контент нередко привязан к конкретным шаблонам и страницам сайта. &
    Структурированные сущности и блоки проще использовать повторно в разных страницах и каналах. &
    Маркетинговые страницы, статьи, проекты и вакансии требуют общей контентной модели, а не набора изолированных шаблонов. \\
    \hline
    \texttt{Preview} и управляемая публикация &
    Базовый \texttt{preview} проще, так как редактирование и рендеринг находятся в одном контуре; публикация тесно связана с runtime \texttt{CMS}. &
    \texttt{Preview} требует отдельной интеграции с витриной, но публикацию можно оформить как явный pipeline между \texttt{CMS} и публичным слоем. &
    Для диплома важен не только просмотр черновика, но и формализованный путь ``черновик -> проверка -> публичная витрина''. \\
    \hline
    Управляемость \texttt{SEO} и метаданных &
    Базовые настройки часто доступны из коробки, но их логика сильнее связана с темой и шаблонами. &
    \texttt{SEO}-поля и метаданные нужно моделировать явно, после чего ими можно управлять как частью контентной сущности. &
    Для маркетингового сайта критично редакторски управляемое \texttt{SEO}: \texttt{canonical}, \texttt{meta} и \texttt{Open Graph} без правки frontend-кода. \\
    \hline
    Удобство для воспроизводимого deployment-контура &
    Стартовый deploy проще как одного приложения, но публичная выдача и администрирование остаются в одном runtime-контуре. &
    Контуры \texttt{CMS} и витрины разводятся, что требует интеграции, но упрощает versioned deployment и изоляцию публикации. &
    Для rebuild-based публикации важнее предсказуемый и воспроизводимый pipeline, чем минимальное число исполняемых компонентов. \\
    \hline
  \end{tabular}
\end{table}

Матрица показывает, что традиционный подход сохраняет преимущество в стартовой простоте
и во встроенности базового \texttt{preview}, однако по критериям, определяющим предметную
область данной ВКР, решающими оказываются другие свойства: независимость frontend-слоя,
сквозная мультиязычность, переиспользование контента, редакторски управляемое \texttt{SEO} и
воспроизводимая публикация. Эти требования могут быть реализованы и поверх традиционной
\texttt{CMS}, но для нее они обычно выступают как надстройки над связанным шаблонным
контуром. Для \texttt{headless CMS} те же свойства являются прямым следствием разделения
контентного и презентационного слоев. Следовательно, выбор \texttt{headless CMS} в рамках
данной работы является не данью технологическому тренду, а рациональным ответом на
критерии, заданные архитектурой корпоративного маркетингового сайта.

\section{Обоснование выбора Strapi и Astro}

\subsection{Strapi как ядро управления контентом}

Выбор \texttt{Strapi} в качестве \texttt{CMS}-ядра обусловлен тем, что система изначально
ориентирована на сценарий \texttt{headless CMS} и предоставляет набор возможностей,
необходимых для рассматриваемой ВКР. Согласно официальной документации \texttt{Strapi 5},
платформа поддерживает моделирование контентных типов, административное управление
содержимым, \texttt{Draft \& Publish}, международную локализацию, \texttt{Preview},
\texttt{Role-Based Access Control} и работу через контентные \texttt{API}~\cite{strapi-cms-docs}.
Для дипломного проекта это означает, что ключевые требования к системе могут быть закрыты
за счет настройки и расширения готового \texttt{CMS}-ядра, а не через разработку
самописной административной панели.

Дополнительным преимуществом \texttt{Strapi} является ориентация на структурированные
контентные сущности. Для предметной области маркетингового сайта это позволяет явно
выделять статьи, проекты, вакансии, отклики, глобальные данные и страницы
с конструктором блоков. Такой подход лучше соответствует цели диплома, чем хранение
контента в виде набора слабо структурированных полей, жестко привязанных к шаблонам.

Для выбранной архитектурной модели также существенно, что \texttt{Strapi} поддерживает
расширяемую \texttt{API}-ориентированную интеграцию, в которой административный слой,
контентная модель и правила доступа могут развиваться независимо от публичной витрины.
Это делает платформу рациональной основой для авторской \texttt{CMS-first} системы без
необходимости разрабатывать \texttt{CMS}-инфраструктуру с нуля.

\subsection{Astro как публичная витрина и monorepo как среда разработки}

\texttt{Astro} выбран в качестве публичной витрины, поскольку его базовая модель хорошо
соответствует контентным сайтам, где значительная часть страниц может быть подготовлена
заранее. Официальная документация \texttt{Astro} указывает, что по умолчанию страницы,
маршруты и endpoint-ы предварительно рендерятся во время сборки как статические страницы,
а переход к рендерингу по запросу выполняется только при необходимости через адаптер и
отключение prerender для соответствующих маршрутов~\cite{astro-on-demand}. Для маркетинговой
платформы это является важным преимуществом: основной контентный слой может публиковаться
как предсобранная витрина, а не зависеть от серверного формирования HTML на каждый запрос.

Еще один значимый аргумент состоит в том, что \texttt{Astro} естественно поддерживает
гибридную модель frontend-разработки. В рассматриваемом проекте это выражается в
использовании \texttt{React}-интеграции и selective hydration через директивы
\texttt{client:*}, что позволяет оставить статическую основу страницы и при этом включать
интерактивность только там, где она действительно требуется. Для маркетингового сайта такой
баланс особенно полезен: информационные и \texttt{SEO}-значимые части страницы остаются
предсказуемыми, а сложные элементы интерфейса можно загружать избирательно.

Практическую обоснованность связки \texttt{Strapi + Astro} подтверждает и официальное
руководство по их совместному использованию, в котором для статического режима прямо
предполагается автоматическая пересборка сайта по \texttt{webhook} после изменения контента~\cite{astro-strapi-guide}.
Для выбранной архитектурной модели это принципиально: публикация должна обновлять
предсобранную витрину автоматически, без ручного вмешательства в frontend-код. Тем самым
схема \texttt{webhook -> rebuild} выступает не частной технической деталью, а логическим
следствием разделения ролей между \texttt{Strapi} как \texttt{CMS}-ядром и \texttt{Astro}
как публичной витриной.

Средой совместной разработки обеих частей проекта выбран монорепозиторий \texttt{Nx + pnpm}.
Официальная документация \texttt{Nx} подчеркивает, что монорепозиторий позволяет хранить
несколько приложений и библиотек в одном репозитории, выполнять атомарные изменения,
обеспечивать единый набор зависимостей и использовать согласованные команды сборки и
тестирования~\cite{nx-monorepos}. Для дипломного проекта это важно по двум причинам.
Во-первых, \texttt{CMS} и витрина развиваются совместно и зависят от общего контракта
данных. Во-вторых, такая организация упрощает генерацию типизированного \texttt{API}-клиента
из \texttt{OpenAPI} и снижает риск рассинхронизации между backend- и frontend-частью.

Следовательно, выбранная связка \texttt{Strapi + Astro} задает архитектурную модель, в
которой управление структурированным контентом, публичная витрина и публикационный контур
разведены по ролям. На основе анализа предметной области, выбранной архитектурной модели и
с учетом целевых редакторских и публикационных сценариев далее формируются требования к
разрабатываемой системе.

\section{Формирование требований к системе}

\subsection{Функциональные требования}

На основе анализа предметной области, выбранной архитектурной модели и с учетом целевых
редакторских и публикационных сценариев к разрабатываемой
системе можно сформулировать следующие функциональные требования:
\begin{itemize}
  \item управление сущностями \texttt{pages}, \texttt{articles}, \texttt{projects},
  \texttt{vacancies} и \texttt{vacancy applications} через \texttt{CMS};
  \item создание маркетинговых страниц на основе \texttt{Dynamic Zone} и набора
  переиспользуемых блоков;
  \item хранение и редактирование глобальных данных сайта в централизованной модели;
  \item поддержка локалей \texttt{ru} и \texttt{en} для публичного контента и основных
  пользовательских текстов;
  \item управление \texttt{SEO}-полями, \texttt{Open Graph}-метаданными и маршрутами;
  \item поддержка \texttt{preview mode} для черновых материалов;
  \item публикация изменений через схему \texttt{webhook -> rebuild};
  \item поддержка ролей и прав доступа для редакторских сценариев;
  \item сохранение и обработка откликов на вакансии с базовой защитой публичной формы.
\end{itemize}

Фактически эти требования отражают два взаимосвязанных пользовательских контура. Первый
контур принадлежит контент-менеджеру и маркетологу: им необходимо создавать и публиковать
материалы без участия разработчика. Второй контур принадлежит внешнему
посетителю: он должен получать структурированные страницы, статьи, кейсы и вакансии в
стабильном и предсказуемом публичном интерфейсе. Следовательно, функциональные требования
должны закрывать одновременно задачи редактора и задачи витрины.

\subsection{Нефункциональные требования}

Нефункциональные требования определяются выбранной архитектурой и масштабом предметной
области. Для рассматриваемой системы существенны следующие свойства:
\begin{itemize}
  \item ориентация на малую или среднюю корпоративную нагрузку без искусственного
  усложнения под высоконагруженный сценарий;
  \item воспроизводимый процесс развертывания с разделением frontend и \texttt{CMS};
  \item базовый уровень безопасности, достаточный для промышленной эксплуатации административных и публичных сценариев;
  \item сопровождаемость за счет переиспользуемых блоков, централизованной модели данных и
  типизированного \texttt{API}-контракта;
  \item приемлемая производительность и индексируемость публичных страниц;
  \item предсказуемость отдачи контентных страниц за счет предварительной сборки витрины.
\end{itemize}

Для публичной витрины дополнительно значима доступность интерфейса; в качестве
нормативного ориентира для проверки таких страниц целесообразно использовать
\texttt{WCAG 2.2}~\cite{wcag22}.

Важным нефункциональным требованием является отделение редакторского контура от публичной
выдачи. Пользовательская страница не должна зависеть от ручного запуска локальных правок
во frontend-коде для типовых редакторских изменений. Аналогично, работа административной
\texttt{CMS} не должна напрямую определять структуру публичного слоя. Для контентных
маркетинговых страниц приоритетной становится устойчивость и воспроизводимость
публикационного процесса, а не максимальная динамичность каждого запроса.

\subsection{Ограничения и нецели проекта}

Чтобы не размывать предмет диплома, необходимо явно зафиксировать границы разрабатываемой
системы. В объем работы не входят:
\begin{itemize}
  \item функции \texttt{e-commerce};
  \item личные кабинеты и пользовательская аутентификация на стороне витрины;
  \item сложные процессы согласования;
  \item совместное редактирование в реальном времени.
\end{itemize}

Одновременно важно подчеркнуть, что за рамками работы не выводятся \texttt{preview mode},
базовые редакторские сценарии, защита публикации и защита публичных форм. Напротив, именно
они формируют прикладной характер системы управления контентом и позволяют рассматривать
разрабатываемую платформу как инженерно завершенное решение для корпоративного
маркетингового сайта.

\section{Выводы по аналитической главе}

Проведенный анализ показывает, что корпоративный маркетинговый сайт следует рассматривать
как контентную систему, а не как набор редко изменяемых страниц. Для такой предметной
области существенны регулярные публикации, мультиязычность, управление метаданными,
сценарий предварительного просмотра и воспроизводимая доставка контента на публичную
витрину. Эти свойства затруднительно обеспечивать последовательно в архитектуре, где
слой управления контентом жестко связан со слоем отображения.

Сопоставление подходов к управлению контентом показывает, что \texttt{headless CMS}
лучше соответствует задачам рассматриваемой предметной области, поскольку позволяет
разделить контентную модель и frontend-представление, использовать предварительную сборку
маркетинговых страниц и выстроить публикационный контур через \texttt{API} и
\texttt{webhook -> rebuild}. В этом контексте выбор \texttt{Strapi} обоснован как выбор
\texttt{CMS}-ядра с поддержкой структурированных сущностей, локализации и редакторских
возможностей, а выбор \texttt{Astro} --- как выбор публичной витрины, ориентированной на
предсобранное представление контента и контролируемую публикацию.

На этой основе в главе сформулированы функциональные и нефункциональные требования к
системе: управление ключевыми контентными сущностями, сборка страниц из блоков,
поддержка \texttt{ru/en}, \texttt{preview mode}, \texttt{SEO}, ролей и прав доступа,
а также воспроизводимого публикационного и эксплуатационного контура. Одновременно
зафиксированы ограничения и нецели, позволяющие удержать инженерную целостность решения
без размывания темы диплома.

Следовательно, во второй главе внимание может быть сосредоточено на проектировании
архитектуры, модели данных, редакторских сценариев и публикационного контура как на
средствах реализации сформированных требований.
